\section{Introduction}
\label{sec:intro}

Simultaneous speech translation (SST) must commit target-language output before
hearing the future of the source stream. Every design choice in the field ---
wait-$k$ schedules \citep{ma2019stacl}, adaptive READ/WRITE policies
\citep{liu2020lowlatency,papi2023alignatt}, long-form streaming architectures
\citep{ouyang2025infinisst} --- is a way of coping with this single constraint:
at commit time, the model knows only a prefix. The words most damaged by this
constraint are the ones that matter most in talks and lectures: rare
terminology, named entities, numbers --- exactly the content a translator cannot
guess from a truncated prefix and cannot revise after committing.

Yet in the settings where SST is most needed --- conference talks, lectures,
screen-shared meetings --- the future is, quite literally, projected on the
wall. Speakers put up a slide and then discuss it, typically for 30--60 seconds;
professional simultaneous interpreters exploit this routinely, requesting decks
in advance and reading the projected slides from the booth. The visual channel
is thus a \emph{legal lookahead}: it can reveal upcoming terminology before it
is spoken. Visual processing is not free, but its cold cost can be paid once
after a slide change, amortized over the slide's dwell window, and kept from
blocking the audio decision loop. We call this \emph{vision-as-anticipation},
and this paper builds the method, the evidence, and the evaluation protocol for
it.

Prior work has approached the ingredients separately but never the combination.
Slide context improves \emph{offline ASR} of conference talks
\citep{sinhamahapatra2025slides}, but transcription has no latency dimension and
cannot exploit the cross-lingual structure of slides. Images help
\emph{simultaneous text MT} on caption datasets \citep{caglayan2020simultaneous,
ive2021exploiting}, establishing the anticipation insight in a toy regime ---
single photographs, no speech, sentence-level decoding. Most recently,
OmniFusion \citep{koneru2025omnifusion} fuses a multimodal foundation model with
a translation LLM for En$\rightarrow$X SimulST, but processes the image
\emph{synchronously on the decoding path} (paying added latency for vision) and
attributes its speed-up to removing the ASR--MT cascade, not to the visual
channel itself. This leaves a narrower unresolved question: whether a
change-triggered visual state that persists across speech chunks improves
SimulST after strong OCR, wrong/stale-slide, causal-timing, and
computation-aware controls.

We decompose the benefit of slide evidence into three mechanisms:
\textbf{M1 anticipation} (slides temporally precede the speech that discusses
them; direction-agnostic), \textbf{M2 recognition support} (source-language
slide terms help the model resolve rare or mis-heard words), and
\textbf{M3 target-form supply} (when the slide language equals the target
language --- common in X$\rightarrow$En, since English is the lingua franca of
slides --- the slide hands the system the target surface form of an upcoming
term, turning translation into recognition plus copying). The decomposition is
testable: slide language is a per-talk stratification variable, and M3 predicts
larger terminology gains and higher slide-string copy rates precisely on talks
whose slides are in the target language.

Our probes show why these controls are necessary. Transcript-input experiments
found large gains from oracle terms and selected VLM-derived terms, but remove
the acoustic ambiguity that semantic evidence is supposed to resolve. A later
speech-input Qwen3-Omni probe found that both the current slide and a different
slide from the same lecture improve over audio-only, while current versus wrong
is statistically indistinguishable ($+0.34$ chrF, $p{=}0.22$). This does not
establish local slide-content use; it leaves domain priming, image presence,
and repeated per-segment visual encoding confounded. The target study therefore
tests a persistent once-per-slide state against strong OCR, text-equivalent,
stale, and cross-talk controls before making a positive visual claim.

Our contributions:
\begin{enumerate}[leftmargin=*]
    \item \textbf{Setting and mechanism analysis.} We formulate vision-as-anticipation
    SST --- a change-triggered semantic state cached across speech chunks --- and
    decompose its benefit into anticipation, recognition support, and
    target-form supply, each separately measurable via slide-language
    stratification (\S\ref{sec:task}).
    \item \textbf{Benchmark suite.} A direction-general evaluation over real
    slides only: mTEDx-V (long-form X$\rightarrow$En, live video-recoverable
    talks, human references, visual-signal stratification), Chinese-LiPS-Long
    (dedicated 1080p slide stream rebuilt on the original session timeline),
    and ACL 60/60 (En$\rightarrow$X control with third-party tagged
    terminology), with a Wikidata-based cross-lingual term protocol
    (\S\ref{sec:benchmark}).
    \item \textbf{Method and evidence.} A non-blocking slide-evidence pipeline
    with one-shot structured reading, persistent caching, and causal retrieval,
    evaluated with content, temporal-alignment, and computation controls that
    separate cold, amortized, and on-path costs
    (\S\ref{sec:method}, \S\ref{sec:results}).
\end{enumerate}
